% Slide 04: Knowledge Grounding Paradigms
\begin{frame}{Knowledge Grounding Paradigms}
  \vspace{-0.1cm}
  \begin{table}[H]
    \centering
    \footnotesize
    \renewcommand{\arraystretch}{1.2}
    \begin{tabularx}{\textwidth}{>{\bfseries}l X l l}
      \toprule
      \rowcolor{tableheader}
      Paradigm & Operating Mechanism & Latency / Compute & Knowledge Update \\
      \midrule
      Fine-Tuning & Updates parametric neural weights & High training cost & Requires retraining \\
      Long-Context & Appends all source text into prompt & Quadratic token cost & Dynamic per query \\
      RAG (Hybrid) & Pairs top-$k$ retrieval with LLM generation & Bounded context cost & Hot-swappable index \\
      \bottomrule
    \end{tabularx}
  \end{table}

  \vspace{0.1cm}
  \begin{itemize}
    \setlength{\itemsep}{0.2cm}
    \item \textbf{Decoupled Architecture:} Parametric model acts as a language synthesizer; non-parametric index acts as the verifiable knowledge store.
    \item \textbf{Efficiency \& Precision:} Mitigates ``Lost-in-the-Middle'' phenomena and context dilution by filtering large corpora into top-$k$ evidence.
    \item \textbf{Lewis et al.\ Finding:} RAG-Sequence achieved 44.5 Exact Match on Natural Questions, outperforming a 50$\times$ larger parametric T5-11B model (36.6).
  \end{itemize}

  \source{Lewis et al. (NeurIPS 2020), Table 1; Liu et al. (2024), TACL.}
\end{frame}
